\chapter{EhBASIC}\label{cha:basic}

The machine speaks EhBASIC~2.22 --- Lee Davison's Enhanced BASIC --- with
this machine's additions: graphics statements that ride the blitter, sound
on four SIDs, floating point on the MATH unit, and an escape hatch to the
shell.

\begin{type}
RUN EHBASIC
\end{type}

It comes up ready: the banner, \texttt{47103 Bytes free} --- more than
twenty-one C64s more than a C64 --- and a
\verb|Ready| prompt --- the traditional \texttt{Memory size ?} question is
gone; the machine knows.

\section{Ten minutes of it}
\begin{type}
RUN "DEMOS.BAS"
\end{type}

\screen{shots/demos.png}{The demo menu. A key picks an entry; each demo
returns here when it ends.}

The menu chains: \verb|RUN "name"| (or \verb|LOAD| inside a running
program) loads another BASIC program and runs it --- that is how one BASIC
program hands over to another, exactly as on a C64.

\section{The machine's own statements}
\begin{type}
GRAPHICS 2          640x480 bitmap over the text
PLOT X,Y,C          LINE X1,Y1,X2,Y2,C
TRI X1,Y1,X2,Y2,X3,Y3,C
PALETTE I,R,G,B     GCLS            GRAPHICS 0
\end{type}

Colours 0--15 belong to the text screen; demos use 16 and up. \verb|GRAPHICS 0|
puts the text mode and its palette back, whatever the program changed.

\section{The * escape}
Any shell command works from BASIC with \verb|*| in front --- \verb|*DIR|,
\verb|*CD EHBASIC|, \verb|*MON|, \verb|*DUMP a note| --- and \verb|*BYE|
leaves BASIC for the shell (\verb|@| is accepted too, a survivor from an
earlier machine). Escape or Ctrl-C stops a running program (and
silences the SIDs); the program's variables survive for \verb|PRINT|-style
post-mortems.

\section{Editing the program in VI}
\verb|*VI| and \verb|*EDIT|, with nothing after them, edit the program that
is in memory:

\begin{type}
10 PRINT "HELLO"
*VI
\end{type}

\noindent BASIC writes the program to a temp file, opens the editor on it,
and reads it back when you leave --- so what you type in the editor is what
you \verb|LIST| afterwards. Give either one a file name and nothing special
happens: \verb|*VI notes.txt| is the ordinary \verb|*| escape, and the
editor edits that file.

The trip out and back is not free. The editors load where the interpreter
lives, so the shell command underneath is \verb|SWAP|, which puts all
64\,KB and the screen aside first and restores them afterwards. And because
the program is written out and read back, \emph{variables do not survive} ---
this is an immediate-mode thing, like \verb|LOAD|. The temp file
(\pth{EDITTMP.BAS}) is left behind, which is occasionally a useful accident.

\begin{aside}
\textbf{Why the first line is blank:} \verb|SAVE| starts its output with a
newline, so the editor opens on an empty line~1 with the program below it.
Harmless --- BASIC ignores it on the way back in --- but it is why the line
count is one more than you expect.
\end{aside}
